% Part: proof-theory
% Chapter: natural-deduction
% Section: quantifiers

\documentclass[../../../include/open-logic-section]{subfiles}

\begin{document}

\olfileid[tr]{pt}{nat}{qua}

\olsection{Düzenli \usetoken{P}{proof} ve Yerine Koyma}

Niceleyicilerin kuralları, \Elim{\lexists} ve~\Intro{\lforall} kuralları
için geçerli olan ``özdeğişken koşulları'' nedeniyle bazı güçlükler doğurur.
Bu güçlüklerin çoğu, bu çıkarımlardaki !!{constant}~$c$'nin hiçbir zaman
yeniden kullanılmaması sağlanarak ele alınabilir. Şu tanımı verelim.

\begin{defn}
Bir doğal tümdengelim !!{proof}{}ı~$\delta$, aşağıdaki koşulları sağlıyorsa
\emph{düzenlidir}:
\begin{enumerate}
  \item hiçbir özdeğişken~$a$, birden çok \Elim{\lexists} ya da
  \Intro{\lforall} çıkarımının özdeğişkeni olarak kullanılmaz ve
  \item $a$, bir \Elim{\lexists} ya da \Intro{\lforall} çıkarımının
  özdeğişkeniyse yalnızca sırasıyla minör öncüle ya da öncüle götüren doğrudan
  alt-!!{proof} içinde geçer.
\end{enumerate}
\end{defn}

\begin{lem}\ollabel{lem:subst-regular} $\delta$ düzenliyse, $!C$ ile
bitiyorsa, $c$ $\delta$ içinde özdeğişken olarak kullanılmayan bir
!!a{constant}{}se ve $t$, $\delta$'nın hiçbir özdeğişkenini içermeyen bir
terimse, $\delta$ içindeki her $c$ geçişinin $t$ ile değiştirilmesiyle ortaya
çıkan $\Subst{\delta}{t}{c}$ düzenli bir !!{proof}{}tır.
$\Subst{\delta}{t}{c}$'nin son-!!{formula}{}ü $\Subst{!C}{t}{c}$'dir.
$!A^x$, $\delta$'nın açık bir varsayımıysa $\Subst{!A}{t}{c}^x$ de
$\Subst{\delta}{t}{c}$'nin açık bir varsayımıdır.
\end{lem}

\begin{proof}
  $\pheight{\delta}$ üzerinde tümevarımla. $\pheight{\delta} = 0$ ise
  $\delta$ yalnızca bir $!A^x$ varsayımından oluşur. Bu durumda
  $\Subst{\delta}{t}{c}$, $\Subst{!A}{t}{c}^x$'tir.

  $\pheight{\delta} > 0$ ise $\delta$'daki son çıkarıma göre durumları ayırırız.
  Önerme bağlaçlarının kurallarına ilişkin durumlar kolaydır ve alıştırma
  olarak bırakılmıştır. $\delta$'nın bir niceleyici çıkarımıyla bittiği
  durumları ele alıyoruz.
  \begin{enumerate}
    \item Son çıkarım $\Intro{\lforall}$ ise $\delta$ şu biçimdedir:
    \[
    \AxiomC{}
    \RightLabel{$\delta'$}
    \DeduceC{$!A(a,c)$}
    \RightLabel{\Intro{\lforall}}
    \UnaryInfC{$\lforall[x][!A(x,c)]$}
    \DisplayProof
    \]
    Tümevarım varsayımı gereği $\Subst{\delta'}{t}{c}$, $!A(a,t)$'nin
    düzenli bir !!{proof}{}ıdır. $a$, $\delta$'nın bir özdeğişkeni olduğundan
    $t$ içinde geçmez. Dolayısıyla $\Subst{\delta}{t}{c}$'deki son çıkarım
    \[
    \AxiomC{}
    \RightLabel{$\Subst{\delta'}{t}{c}$}
    \DeduceC{$!A(a,t)$}
    \RightLabel{\Intro{\lforall}}
    \UnaryInfC{$\lforall[x][!A(x,t)]$}
    \DisplayProof
    \]
    doğrudur: $t$, $a$'yı içermediğinden ne sonuç ne de herhangi bir açık
    varsayım $a$'yı içerir.
    \item Son çıkarım $\Elim{\lforall}$ ise $\delta$ şu biçimdedir:
    \[
    \AxiomC{}
    \RightLabel{$\delta'$}
    \DeduceC{$\lforall[x][!A(x,c)]$}
    \RightLabel{\Elim{\lforall}}
    \UnaryInfC{$!A(s(c),c)$}
    \DisplayProof
    \]
    Tümevarım varsayımı gereği $\Subst{\delta'}{t}{c}$,
    $\lforall[x][!A(x,t)]$'nin düzenli bir !!{proof}{}ıdır. (Sonuçtaki $s(c)$
    terimi $c$'yi içerebilir de içermeyebilir de.) $\Subst{\delta}{t}{c}$'deki
    son çıkarım
    \[
    \AxiomC{}
    \RightLabel{$\Subst{\delta'}{t}{c}$}
    \DeduceC{$\lforall[x][!A(x,t)]$}
    \RightLabel{\Elim{\lforall}}
    \UnaryInfC{$!A(s(t),t)$}
    \DisplayProof
    \]
    doğrudur ve $!A(s(t),t) \ident \Subst{!A(s(c),c)}{t}{c}$ olur.
    \Intro{\lexists} durumu benzerdir.
  \item Son çıkarım $\Elim{\lexists}$ ise $\delta$ şu biçimdedir:
    \[
    \AxiomC{}
    \RightLabel{$\delta_1'$}
    \DeduceC{$\lexists[x][!A(x,c)]$}
    \AxiomC{$\Discharge{!A(a,c)}{x}$}
    \RightLabel{$\delta_2'$}
    \DeduceC{$!C$}
    \DischargeRule{\Elim{\lexists}}{x}
    \BinaryInfC{$!C$}
    \DisplayProof
    \]
    Tümevarım varsayımı gereği $\Subst{\delta_1'}{t}{c}$ ve
    $\Subst{\delta_2'}{t}{c}$, sırasıyla $\lexists[x][!A(x,t)]$'nin ve açık
    $!A(a,t)$ varsayımından $\Subst{!C}{t}{c}$'nin düzenli !!{proof}s{}ıdır.
    $a$, $\delta$'nın bir özdeğişkeni olduğundan $t$ içinde geçmez.
    Dolayısıyla $\Subst{\delta}{t}{c}$'deki son çıkarım
    \[
    \AxiomC{}
    \RightLabel{$\Subst{\delta_1'}{t}{c}$}
    \DeduceC{$\lexists[x][!A(x,t)]$}
    \AxiomC{$\Discharge{!A(a,t)}{x}$}
    \RightLabel{$\Subst{\delta_2'}{t}{c}$}
    \DeduceC{$\Subst{!C}{t}{c}$}
    \DischargeRule{\Elim{\lexists}}{x}
    \BinaryInfC{$\Subst{!C}{t}{c}$}
    \DisplayProof
    \]
    doğrudur: $\Subst{!A(x,t)}{a}{x} \ident
    \Subst{!A(a,c)}{t}{c}$ olur ve ne sonuç~$\Subst{!C}{t}{c}$ ne de
    herhangi bir açık varsayım~$a$'yı içerir.
  \end{enumerate}
\end{proof}

\begin{prop}\ollabel{prop:regular}
$\delta$, \Log{N1c} ya da \Log{N1i}'de bir !!a{proof} ise düzenli bir
!!{proof}~$\delta'$ vardır; bu ispat $\delta$ ile aynı yapıya (özellikle aynı
!!{height}~değerine) sahiptir ve ondan yalnızca özdeğişkenlerin değiştirilmesi
bakımından ayrılır.
\end{prop}

\begin{proof}
Yalnızca bu ispat için, $\delta$'nın bir özdeğişken çıkarımına, özdeğişkeni
başka bir çıkarımın özdeğişkeni değilse ve $\delta$ içinde o çıkarımın
doğrudan alt-!!{proof}{}ı dışında hiçbir yerde geçmiyorsa \emph{temiz}, aksi
hâlde \emph{kirli} diyelim. $n$, $\delta$'daki kirli özdeğişken çıkarımlarının
sayısı olsun. $n$ üzerinde tümevarımla $\delta$'nın gerekli düzenli
$\delta'$ !!{proof}{}ına dönüştürülebileceğini gösteriyoruz. $n=0$ ise
$\delta$'daki bütün özdeğişken çıkarımları temizdir; dolayısıyla $\delta$
düzenlidir ve ispatlanacak bir şey yoktur. Aksi hâlde en üstteki özdeğişken
çıkarımlarından birini, örneğin \Intro{\lforall}'ı seçin. $\delta$'nın bu
çıkarımla biten alt-!!{proof}{}ı~$\delta_1$ şu biçimdedir:
    \[
    \AxiomC{}
    \RightLabel{$\delta_2$}
    \DeduceC{$!A(c)$}
    \RightLabel{\Intro{\lforall}}
    \UnaryInfC{$\lforall[x][!A(x)]$}
    \DisplayProof
    \]
$c$ bir özdeğişken olduğundan $\lforall[x][!A(x)]$ içinde ya da açık bir
varsayımda geçmez. $\delta_2$ ispatı düzenlidir; çünkü hiçbir özdeğişken
çıkarımı içermez. $a$, $\delta$ içinde geçmeyen bir !!a{constant} olsun.
\olref{lem:subst-regular} gereği $\Subst{\delta_2}{a}{c}$, $!A(a)$'nın
düzenli bir ispatıdır. Özdeğişken koşulu gereği $\delta_2$'nin hiçbir açık
varsayımı $c$'yi içermez; dolayısıyla $\Subst{\delta_2}{a}{c}$'nin hiçbir açık
varsayımı da $a$'yı içermez. Böylece~\Intro{\lforall} uygulayabiliriz.
$\delta$ içindeki $\delta_1$'i şu $\delta_1'$ ispatıyla değiştirirsek
    \[
    \AxiomC{}
    \RightLabel{$\Subst{\delta_2}{a}{c}$}
    \DeduceC{$!A(a)$}
    \RightLabel{\Intro{\lforall}}
    \UnaryInfC{$\lforall[x][!A(x)]$}
    \DisplayProof
    \]
kirli bir özdeğişken çıkarımını temiz bir çıkarımla değiştirmiş oluruz ve tek
fark, $c$ özdeğişkeninin~$a$ ile değiştirilmesidir. Ortaya çıkan ispatta
$n-1$ kirli özdeğişken çıkarımı vardır; sonuç tümevarım varsayımından çıkar.
\end{proof}

\begin{prob}
\olref{prop:regular} ispatında en üstteki özdeğişken çıkarımının
\Elim{\lexists} olduğu durumu betimleyin.
\end{prob}

Az önce ispatlanan önermeye açıkça başvurmayacağız; ancak ele alınan bütün
!!{proof}s{}ın düzenli olduğunu varsayacağız---düzenli değillerse
özdeğişkenlerini değiştirerek onları her zaman düzenli kılabiliriz.

\Olref{lem:subst-regular} ve \olref{prop:regular}, \Log{N2c} ve~\Log{N2i}
için de geçerlidir. İspatları alıştırma olarak bırakıyoruz.

\end{document}
